\documentclass[12pt,a4paper]{article}
\usepackage[utf8]{inputenc}
\usepackage{geometry}
\usepackage{hyperref}

\geometry{margin=1in}

\title{\textbf{Literature Review: Reliability-Aware Charging Planning and AI-Based EV Infrastructure}}
\author{Jia Chinnu Abraham}
\date{\today}

\begin{document}

\maketitle

\section{Introduction}
This document reviews key research papers focusing on smart electric vehicle (EV) charging infrastructure, real-time availability tracking, port reservation systems, and reliability-aware route planning for Project Phase 1.

\section{Reviewed Papers}

\subsection{1. TrustCharge: Reliability-Aware Long-Trip Charging Planning and Port Reservation for Battery Electric Vehicles Using Live Open Data}
\begin{itemize}
    \item \textbf{Focus:} Centralized EV charging reliability, live open data integration, and smart slot reservation mechanisms.
    \item \textbf{Key Contributions:} Proposes a reliability-aware framework leveraging real-time open data to minimize charging anxiety and optimize long-trip scheduling for battery electric vehicles.
\end{itemize}

\subsection{2. AI-Based Smart Electric Vehicle Charging: A Review of Load Balancing, Reservation, Reliability, and Waiting Time Prediction}
\begin{itemize}
    \item \textbf{Focus:} Comprehensive overview of artificial intelligence techniques applied to EV charging load balancing, reservation scheduling, and queue wait-time estimation.
    \item \textbf{Key Contributions:} Analyzes various optimization algorithms and machine learning models used to predict station congestion and distribute grid loads efficiently.
\end{itemize}

\end{document}
